%! Matrices and Column Spaces — Unit 1, Lesson 1.3 — Lesson Plan
\documentclass[10pt]{article}
\usepackage{linalg-article}
\usepackage{linalg-boxes}
\usepackage{graphicx}   % -article does NOT load graphicx; needed for thumbnails/screenshots
\graphicspath{{images/}}
\newcommand{\TallMath}[1]{$\displaystyle #1\rule[-1.4em]{0pt}{3.2em}$}
\newcommand{\vv}{\mathbf{v}}
\newcommand{\ww}{\mathbf{w}}
\newcommand{\xx}{\mathbf{x}}
\newcommand{\bb}{\mathbf{b}}
\newcommand{\UnitNumberName}{Unit 1: Vectors and Matrices \quad}
\newcommand{\LessonNumberName}{Lesson 1.3: Matrices and Column Spaces}

\begin{document}
\begin{center}
    {\Large\bfseries \CourseName: \SchoolYear \\
    {\small\bfseries \UnitNumberName \LessonNumberName}}
\end{center}

\begin{tcolorbox}[colback=blush, colframe=burgundy, boxrule=0.9pt, arc=2mm,
    left=2mm, right=2mm, top=1.5mm, bottom=1.5mm]
    \textbf{Primary Objective:} Students can read a matrix as vectors standing side by side
    (its \emph{columns}), compute $A\xx$ as a \emph{linear combination of those columns}
    $x_1\mathbf{a}_1+x_2\mathbf{a}_2$, and describe the \emph{column space} $C(A)$ --- every
    target $\bb$ the columns can build. They can justify \emph{geometrically} why two
    independent columns fill the whole plane while two parallel columns fill only a line, and
    decide whether a given $\bb$ is reachable.
\end{tcolorbox}

\renewcommand{\arraystretch}{1.4}
\begin{skillbox}[Priority Ideas \& Skills]{goldbox}
\begin{tabularx}{\linewidth}{@{}X|X@{}}
\textbf{Priority Ideas \& Skills} & \textbf{Key Understandings (the ``why'')} \\[2pt]
\hline
\vspace{-6pt}
\begin{itemize}[leftmargin=1.1em, nosep, after=\vspace{2pt}]
  \item Read a matrix as its \emph{columns} $\mathbf{a}_1,\mathbf{a}_2$.
  \item Compute $A\xx = x_1\mathbf{a}_1 + x_2\mathbf{a}_2$.
  \item Describe the column space $C(A)$ as all combinations of the columns.
  \item Decide if $C(A)$ is a \emph{line} or the whole \emph{plane}, and whether a target $\bb$ is reachable.
\end{itemize}
&
\vspace{-6pt}
\begin{itemize}[leftmargin=1.1em, nosep, after=\vspace{2pt}]
  \item A matrix is just vectors lined up as columns --- nothing new, repackaged.
  \item $A\xx$ \emph{mixes} the columns with the weights in $\xx$ --- the same linear combination from Lesson~1.1.
  \item The column space is the \emph{span of the columns}: every output the matrix can produce.
  \item Independent columns fill the plane; parallel columns fill only a line --- and $A\xx=\bb$ is solvable exactly when $\bb$ sits in $C(A)$.
\end{itemize}
\end{tabularx}
\end{skillbox}

\begin{skillbox}[{Vocabulary, Concepts, \& Theorems}]{greenbox}
\renewcommand{\arraystretch}{1.3}
\begin{tabularx}{\linewidth}{@{}l X@{}}
\textbf{Matrix} & A rectangular array of numbers; read it as vectors standing side by side --- its \emph{columns}. \\
\textbf{Columns of $A$} & The vectors $\mathbf{a}_1,\mathbf{a}_2$ that make up $A=\begin{bsmallmatrix} \mathbf{a}_1 & \mathbf{a}_2 \end{bsmallmatrix}$. \\
\textbf{Matrix--vector product} & \TallMath{A\xx = x_1\mathbf{a}_1 + x_2\mathbf{a}_2}; a linear combination of the columns, weighted by $\xx$. \\
\textbf{Column space $C(A)$} & All vectors $A\xx$ --- every combination of the columns; the \emph{span} of the columns. \\
\textbf{Independent / dependent columns} & Independent (not parallel) columns span the whole plane; dependent (parallel) columns span only a line. \\
\end{tabularx}
\end{skillbox}

\begin{fixedskillbox}[Activate Prior Knowledge \& Spiral Review]{blush}
    The Warm-Up rehearses the Lesson~1.1 skills this lesson repackages as a matrix:
    \textbf{(1)} forming a \emph{linear combination} $c\vv+d\ww$ (which becomes $A\xx$),
    \textbf{(2)} deciding whether one vector is a \emph{multiple} of another (parallel), and
    \textbf{(3)} recalling that two \emph{non-parallel} vectors span the whole plane while two
    parallel ones span a line. Debrief item~3 aloud --- ``line vs.\ plane'' returns immediately
    as the shape of the column space.
\end{fixedskillbox}

\begin{skillbox}[Hook]{blush}
    Put two columns on the board as a matrix,
    $A=\begin{bsmallmatrix} 1 & 3 \\ 2 & 1 \end{bsmallmatrix}$, and ask:
    \textbf{``If a matrix is really just two vectors standing side by side, what does it mean to
    multiply it by $\xx=\begin{bsmallmatrix} 2\\1 \end{bsmallmatrix}$?''} Reveal the ``columns''
    rule: $A\xx = 2\begin{bsmallmatrix} 1\\2 \end{bsmallmatrix} + 1\begin{bsmallmatrix} 3\\1 \end{bsmallmatrix}
    = \begin{bsmallmatrix} 5\\5 \end{bsmallmatrix}$ --- a linear combination of the columns.
    Then push the big question: \emph{as $\xx$ ranges over everything, which targets $\bb$ can we
    hit?} That collection is the \textbf{column space}. This frames the lesson: a matrix packages
    vectors, and its column space is everything they can build.
\end{skillbox}

\begin{skillbox}[Lesson]{blush}
\begin{multicols}{2}
\textbf{1. A matrix is vectors in columns.} Write $A=\begin{bsmallmatrix} 1 & 3 \\ 2 & 1 \end{bsmallmatrix}$
as columns $\mathbf{a}_1=\begin{bsmallmatrix} 1\\2 \end{bsmallmatrix}$,
$\mathbf{a}_2=\begin{bsmallmatrix} 3\\1 \end{bsmallmatrix}$. A matrix is not a new object --- it is
vectors lined up.

\textbf{2. $A\xx$ combines the columns.} Define
$A\xx = x_1\mathbf{a}_1 + x_2\mathbf{a}_2$. With $\xx=\begin{bsmallmatrix} 2\\1 \end{bsmallmatrix}$,
$A\xx=2\begin{bsmallmatrix} 1\\2 \end{bsmallmatrix}+1\begin{bsmallmatrix} 3\\1 \end{bsmallmatrix}
=\begin{bsmallmatrix} 5\\5 \end{bsmallmatrix}$. This is the Lesson~1.1 combination, now written with a matrix.

\columnbreak

\textbf{3. The column space.} $C(A)$ is \emph{all} vectors $A\xx$ --- every combination of the
columns, i.e.\ their \textbf{span}. For $A$ above the two columns are not parallel, so $C(A)$ is
the \emph{whole plane}: every $\bb$ is reachable.

\textbf{4. Line vs.\ plane; is $\bb$ reachable?} Contrast $B=\begin{bsmallmatrix} 1 & 2 \\ 2 & 4 \end{bsmallmatrix}$:
column~2 is $2\times$ column~1 (parallel), so $C(B)$ is only the \emph{line} through
$\begin{bsmallmatrix} 1\\2 \end{bsmallmatrix}$. Then $\begin{bsmallmatrix} 3\\6 \end{bsmallmatrix}$ is
reachable but $\begin{bsmallmatrix} 1\\0 \end{bsmallmatrix}$ is not --- solving $A\xx=\bb$ asks
whether $\bb$ lies in $C(A)$ (previews Unit~2).
\end{multicols}
\end{skillbox}

\begin{skillbox}[Active Monitoring]{redbox}
Circulate during the notes practice and Tier work. Watch for and correct:
\begin{itemize}[leftmargin=1.2em, nosep]
  \item computing $A\xx$ row-by-row and losing the \emph{columns} picture (miss the linear combination);
  \item mixing up which weight scales which column ($x_1$ with $\mathbf{a}_1$);
  \item thinking every matrix's column space is the whole plane --- missing the parallel-columns ``line'' case;
  \item saying $\bb$ is reachable without checking whether it lies on the span of the columns.
\end{itemize}
Cold-call prompts: ``What are the two columns of this matrix?'' ``$A\xx$ is a combination of
\emph{what}?'' ``Line or plane --- how do you know?'' ``Is $\bb$ in the column space?''
\end{skillbox}

\begin{skillbox}[Group Work \& Differentiation]{redbox}
\textbf{Building Orders from Two Boxes} activity (tiered; mirrors the notes). The columns of $A$
are two gift boxes $\begin{bsmallmatrix} \text{mugs} \\ \text{shirts} \end{bsmallmatrix}$, and
$A\xx$ is the total order built from $x_1$ of one box and $x_2$ of the other. The column space is
every order the two boxes can fill.
\begin{multicols}{3}
\textbf{Tier R --- Remediate.} Identify the columns of $A$ and compute $A\xx$ for given $\xx$ ---
mixing the two boxes into a total order.

\columnbreak
\textbf{Tier A --- Approaching.} Solve $A\xx=\bb$ for a target order, and decide whether the two
independent boxes can fill \emph{any} target (whole plane).

\columnbreak
\textbf{Tier E --- Extension.} A shop with two \emph{parallel} boxes: show its column space is a
line, find an order it can never fill, and justify why using the columns.
\end{multicols}
\end{skillbox}

\begin{skillbox}[Individual Work \& Assessment]{redbox}
\textbf{Exit Ticket} (independent, no notes): name the columns of a matrix and compute an $A\xx$
as a combination of columns; then a \textbf{reachability/justification check} --- decide whether a
given $\bb$ lies in the column space of a matrix with parallel columns and explain how you know.
Collect and sort into ``got it / arithmetic slip / concept gap (columns picture, or line-vs-plane)''
to decide whether 1.4 opens with a quick re-teach.
\end{skillbox}

\begin{skillbox}[Reinforcement \& Extension]{goldbox}
\textbf{Homework:} mixed practice --- naming columns, computing $A\xx$ as column combinations,
classifying a column space as a line or the whole plane, and testing whether a target $\bb$ is
reachable. \textbf{Extension:} a \emph{production} context --- two product bundles as columns and
which output totals they can build, extending the unit's data-as-a-table theme. \textbf{Preview:}
Lesson~1.4, \emph{Matrix Multiplication and $A=CR$} --- building each output as a combination of a
few independent columns.
\end{skillbox}
% ── Teacher notes (migrated from the component keys) ──
% One per component, titled for it. They live here rather than in the _key
% files so that every component is the same length blank and keyed.

\begin{teachernote}[Warm-Up]
Every item previews today's matrix. Item~1 \emph{is} $A\xx$: a linear combination of two columns
(1b is exactly $A\begin{bsmallmatrix} 1\\2 \end{bsmallmatrix}$ for $A=\begin{bsmallmatrix} 3&1\\1&2 \end{bsmallmatrix}$).
Item~2 is the parallel test that decides whether a column space is a line. Item~3 is the ``line
vs.\ plane'' fact that becomes the shape of the column space.
\end{teachernote}

\begin{teachernote}[Group Activity]
Tier~A item~3 is the central ``interpret'' moment: independent columns $\Rightarrow$ column space
is the whole plane $\Rightarrow$ every target is reachable. Tier~E flips it: parallel columns
$\Rightarrow$ the column space collapses to a line, so most targets are \emph{unreachable} ---
which is exactly what ``$A\xx=\bb$ has no solution'' will mean in Unit~2. Accept any correct
reasoning in Tier~E item~3 (component equations or the ``$2\times$'' pattern).
\end{teachernote}

\begin{teachernote}[Exit Ticket]
Item~2 checks the central idea: $\bb$ is reachable exactly when it lies in the column space, and
parallel columns collapse that space to a line. Full credit needs both the check (not a multiple
of $\begin{bsmallmatrix} 1\\2 \end{bsmallmatrix}$) \emph{and} the meaning (no solution). A common
slip in item~1 is crossing components --- look for $2\begin{bsmallmatrix} 1\\2 \end{bsmallmatrix}+1\begin{bsmallmatrix} 4\\3 \end{bsmallmatrix}$, not a row-by-row mash.
\end{teachernote}

\end{document}
